% Lösungen Einheit 1 von 3 – Ganzrationale Funktionen (zusammenbau v0.1)
\einheitenkopf{Einheit 1 von 3 · Ganzrationale Funktionen}

\erg{10}{a) $-0{,}1x^3$}
\erg{11}{a) Grad ungerade, Koeffizient negativ \quad b) für $x \to +\infty$: $f(x) \to +\infty$; für $x \to -\infty$: $f(x) \to +\infty$ \quad c) für $x \to +\infty$: $f(x) \to -\infty$; für $x \to -\infty$: $f(x) \to -\infty$ \quad d) für $x \to +\infty$: $f(x) \to -\infty$; für $x \to -\infty$: $f(x) \to +\infty$ \quad e) für $x \to +\infty$: $f'(x) \to +\infty$; für $x \to -\infty$: $f'(x) \to +\infty$ \quad f) alle Exponenten gerade, das Absolutglied zählt als gerade Potenz: achsensymmetrisch zur $y$-Achse \quad g) nur gerade Exponenten (das Absolutglied zählt als gerade Potenz): achsensymmetrisch zur $y$-Achse; gerader Grad, negativer Leitkoeffizient: für $x \to +\infty$: $f(x) \to -\infty$; für $x \to -\infty$: $f(x) \to -\infty$}
\erg{12}{a) Fehler: bei ungeradem Grad laufen die Seiten entgegengesetzt. Richtig: für $x \to +\infty$: $f(x) \to +\infty$; für $x \to -\infty$: $f(x) \to -\infty$ \quad b) $x^4$ wächst schneller als $50x^3$ (der Quotient $\frac{x}{50}$ wird beliebig groß); der Summand mit der kleineren Potenz spielt dann keine Rolle mehr. \quad c) $f(x) = -0{,}25x^4 + 2x^2$}
